\section{Experiments}
\label{sec:exp}

\paragraph{Data and cohorts.}
OLLB: the \texttt{contents} table of the Open LLM Leaderboard v2 (pinned commit \texttt{9c09a7ca}), with its six official normalised task scores (IFEval, BBH, MATH Lvl~5, GPQA, MuSR, MMLU-Pro), 4{,}496 models after removing 79 duplicate names and one flagged model (submissions 2024-06 to 2025-03).
HELM MMLU: release v1.13.0, 79 models, exact-match accuracy on 57 subjects.
Both are frozen snapshots; all ranks and costs refer to them.
A \emph{target} is any model of rank $\ge2$ under the mean rule: 4{,}495 on OLLB and 78 on HELM.
Score rules run on every target.
The ordinal solvers are slower, so Borda and Condorcet run on a rank-stratified sample of 136 OLLB targets (ranks 2--21 and log-uniform tail draws) with stratum weights that make weighted shares and medians estimate the full-leaderboard quantity; HELM runs every rule on every target.
We report by rank group of the target under the mean rule (2--10, 11--100, $>$100).
Caps (Section~\ref{sec:method}): ceiling (100), frontier${}+2$ (best rival${}+2$ points), $+5/10/20$ points over the target's own score, and the empirical 95th and 99th percentiles of fine-tune-minus-base gains (OLLB, 540 base/fine-tune pairs).
For HELM there are no such pairs, so empirical caps use the OLLB MMLU-Pro gain percentile for every subject, a proxy marked $^\dagger$ throughout.
Intervals are 95\% percentile bootstraps over targets within rank strata; they describe these snapshots, not a population of leaderboards.
OLLB values are normalised points and HELM values percentage points, so costs are compared across rules, never across datasets.
The solvers are exact (closed forms, convex or greedy solvers, and HiGHS integer programs), and every solver status is either recorded as optimal or cross-checked against an independent feasibility test (Appendix~\ref{app:repro}).

\paragraph{Headline: how many targets can reach first place at all.}
Figure~\ref{fig:feas} shows the share of targets for which first place is attainable when no task may be raised by more than $h$ points, i.e.\ the distribution of the smallest uniform lift $d^\ast$ (a target is feasible under a gain cap $h$ exactly when $d^\ast\le h$, which we verified on every row).
Table~\ref{tab:feas} adds the empirical cap and the rank groups.
\begin{figure}[!tbp]
\centering
\includegraphics[width=0.8\linewidth]{../figures/F1_feasibility_curve.pdf}
\caption{Share of targets for which first place is attainable when no task may be raised by more than $h$ points (weighted empirical CDF of the minimal uniform lift $d^\ast$; vertical lines at $h=5,10,20$). Top: all targets of rank $\ge2$ (OLLB 4{,}495; HELM 78); bottom: the 29 targets ranked 2--30. Borda and Condorcet on OLLB: 136-target stratified sample, stratum-weighted.}
\label{fig:feas}
\end{figure}
\begin{table}[!tbp]
\centering\footnotesize
\caption{Feasibility of first place by rule and the target's rank under the mean (2--10, 11--100, $>$100): share of targets for which first place is attainable with at most 10 points per task, and under the empirical 95th-percentile gain cap; and the minimal uniform lift $d^\ast$ in points, median [95\% bootstrap interval, omitted where narrower than 1 point]. OLLB: 4{,}495 targets (score rules), 136-target stratified sample (Borda, Condorcet); HELM: 78 targets, none ranked beyond 79. $^\dagger$HELM empirical caps are a proxy.}
\label{tab:feas}
\begingroup\input{../figures/tables/T1_feasibility.tex}\endgroup
\end{table}
At a 10-point cap the shares of all OLLB targets are 3\% (mean), 5\% (median), 10\% (minimum), 3\% (geometric mean), 29\% (Borda) and 3\% (Condorcet winner); on HELM they are 37\%, 51\%, 18\%, 31\%, 31\% and 49\%.
Almost all models far from the top cannot reach first place: of the OLLB targets ranked beyond 100 at most 28\% (Borda; 9\% for the minimum) and under the empirical 95\% cap at most 25\% can (Table~\ref{tab:feas}).
For the 29 best-placed targets (ranks 2--30) the shares are high, 76--100\% on OLLB and 41--100\% on HELM, and the minimum rule is lowest on both.
The rules change \emph{who} can win, not only the price: at a 20-point cap the minimum, Borda and Condorcet rules are attainable for 46--64\% of OLLB targets against 14--21\% for the mean, median and geometric mean, whereas on HELM the minimum rule is the hardest at 5, 10 and 20 points (9\%, 18\% and 44\%).
The empirical cap is stricter: under the 95th-percentile cap only the targets near the top qualify (OLLB ranks 2--10: 100\% for mean, median, geometric mean and Borda, 44--56\% for minimum and Condorcet; ranks 11--100: 22--100\%), and on HELM, with the proxy cap, at most 1\% of targets of rank 11--100 for any rule.
The smallest uniform lift grows with rank under every rule (OLLB median $d^\ast$ of 0--8 points for ranks 2--10, 5--9 for ranks 11--100, 18--36 beyond rank 100; HELM 11--23 points for ranks 11--100).

\begin{figure}[!tbp]
\centering
\includegraphics[width=0.74\linewidth]{../figures/F2_cap_sweep.pdf}
\caption{Cap sweep. Top: share of targets (rank $\ge2$) for which first place is attainable, by rule and cap scenario (caps are not nested; $^\dagger$ proxy). Bottom: median $K^\ast$ on the common feasible cohort of that cap (tick label: cohort size / targets run); points are shown for cohorts of at least 5 targets, bars: 95\% bootstrap intervals. Cohorts differ between caps: compare rules within one cap.}
\label{fig:cap}
\end{figure}

\paragraph{Cost on the common feasible cohort.}
Costs are compared only among targets feasible under every listed rule at that cap (Table~\ref{tab:cost}).
\begin{table}[!tbp]
\centering\footnotesize
\caption{Cost of first place on the common feasible cohort at a cap, median [95\% bootstrap interval] of tasks touched $K^\ast$ and total points $L_1^\ast$. The cohort and the targets run are in each block header. OLLB values are normalised points, HELM values percentage points, so they are not comparable across datasets; cohorts differ between blocks, so compare rules within a block. The ceiling and frontier${}+2$ caps leave every target feasible; cohorts at attainable caps are small and are in Appendix~\ref{app:tables}.}
\label{tab:cost}
\begingroup\input{../figures/tables/T2_common_cohort_cost.tex}\endgroup
\end{table}
At the ceiling cap on OLLB the minimum rule needs the fewest points (25.3 against 100.6 for the mean), then the median (48.1), Borda (65.6), Condorcet (78.1) and the geometric mean (78.8).
On HELM the minimum (98) and median (111) rules need the fewest points and the mean (661) and geometric mean (633) the most, with Condorcet (213) and Borda (500) in between; relative to the mean's $L_1^\ast$, the median rule's median is 0.15 (95\% interval 0.14--0.19) and the minimum's 0.14 on HELM, and 0.60 and 0.25 on OLLB (\texttt{rel\_L1\_vs\_mean}).
The number of tasks tells a different story: Borda needs the most tasks (OLLB 5 of 6; HELM 46 of 57), the minimum rule the fewest on HELM (11), and the mean rule's $K^\ast$ depends strongly on the cap, because a binding cap forces the total to be spread over more tasks (HELM: 16 at the ceiling, 31 at frontier${}+2$).
So the rule that is cheapest in points is not the one that touches the fewest tasks, and the ranking of rules depends on which investment measure is used; on both datasets the ordinal rules need many tasks but not many points relative to the mean.

\paragraph{Sensitivity to the cap, rank, number of tasks and scale.}
Figure~\ref{fig:cap} sweeps all cap scenarios.
Ceiling and frontier caps leave every target feasible, while at $+5$ points 0--2\% of OLLB targets are feasible and at $+20$ points 14--64\%, so the feasibility ordering of the rules shown above is what changes with the cap; we do not rank costs across rules at attainable caps, where cohorts are small (3 to 91 targets) and differ.
Cost rises with the target's rank for every rule overall, though not strictly at every rank (Appendix Figure~\ref{fig:rank}).
On HELM random subsets of $n$ subjects (100 draws per $n$) show that Borda and Condorcet need more tasks as $n$ grows, so part of the Borda cost above reflects $n=57$ (Appendix Figure~\ref{fig:subset}).
On OLLB, replacing the official normalised scores with raw accuracies, per-task $z$-scores or rank percentiles changes median $K^\ast$ for targets ranked 3--10 by at most one task, but for targets ranked beyond 10 by up to four: at the ceiling the mean rule needs 2 tasks on the official scale and 6 on rank percentiles, and the minimum rule 3 against 5 (raw) and 6 ($z$-score) (Appendix Figure~\ref{fig:scale}). We therefore state OLLB conclusions for the official scale and treat the others as sensitivity.


\paragraph{Robustness and theory check.}
Changing the strictness margin $\varepsilon$ from $10^{-3}$ to $10^{-2}$ or $10^{-1}$ changes the median $L_1^\ast$ by at most 12\%; $K^\ast$ changes for at most 13\% of targets at the ceiling cap (Table~\ref{tab:eps}).
Removing merged models leaves median $K^\ast$ at the ceiling and frontier caps unchanged; keeping only the 500 top-ranked models leaves it unchanged for targets ranked 3--10 and lowers it by one to two tasks for lower-ranked targets (Table~\ref{tab:slices}).
An evaluation-noise Monte Carlo and a task bootstrap on 10 targets (100--200 draws) move the feasible share of the minimum rule from the observed 70\% to 49--70\% and leave median $K^\ast$ within one task of the observed value on OLLB and within up to eight tasks (Borda) on HELM (Table~\ref{tab:noise}); these are conditional sensitivity analyses on few targets, not estimates of leaderboard sampling variation.
The closed forms (mean, median for odd $n$, minimum) and the uniform-lift identity agree with the solvers on all 31{,}080 instances (Table~\ref{tab:theory}).
